\documentclass{situationist-preprint}

\title{Interest-Relative Fundamentality}
\msid{MET-10-b}
\msfamily{MET-10}
\msdiscipline{Metaphysics}
\mstype{Consequence}
\mskeywords{fundamentality; levels; special sciences; screening off; grounding; infinite descent}

\begin{document}
\maketitle

\begin{abstract}
MET-10-a argues that no level of reality is fundamental: every level is grounded in a finer one, without end. Yet the sciences routinely treat a level as fundamental, as chemistry treats atoms and population genetics treats genes. We derive from the descent a relative notion: a level is fundamental for an inquiry when the inquiry's explananda are grounded in it and finer facts make no further difference to them. Fundamentality in this sense is relative to the questions asked, it is consistent with the absence of an absolutely fundamental level, and it is the only kind of fundamentality a level can have. A science that treats a level as fundamental commits no metaphysical error.
\end{abstract}

\section{The relative notion}

MET-10-a defines a level as a plurality of non-overlapping proper parts of the cosmos that together compose it (Definition~3.1), and a level as fundamental if no fact at it is partly grounded in a fact at another level (Definition~3.3). Theorem~4.3 of MET-10-a shows that no level is fundamental in that sense. The sciences nonetheless speak of fundamental entities: the atoms of chemistry, the cells of physiology, the agents of economics. The descent tells us what such talk can mean.

\begin{definition}\label{d:rel}
Let $I$ be an inquiry with a class $Q$ of explananda, facts at some level or about objects at it. A level $L$ is \emph{fundamental for $I$} if (i) every member of $Q$ has a full ground consisting of facts at $L$ or at coarser levels, and (ii) any two possible situations that agree in their facts at $L$ agree on every member of $Q$, however they differ at finer levels.
\end{definition}

Clause (ii) is a screening-off condition: given the facts at $L$, the finer facts make no further difference to what $I$ asks about. It does not say that the finer facts ground nothing at $L$; by Principle~P1 of MET-10-a, they ground all of it. It says that the variation the finer facts could introduce, holding $L$ fixed, does not reach the explananda.

\begin{proposition}\label{p:rel}
Fundamentality for an inquiry is consistent with the descent of MET-10-a, and is the only kind of fundamentality a level has.
\end{proposition}

\begin{proof}
Definition~\ref{d:rel} imposes no condition on the grounds of facts at $L$, so it is satisfied by a level all of whose facts are grounded in a finer level, as MET-10-a requires of every level. Its clauses are satisfiable together: in the model of MET-10-a (Proposition~4.5), a cube decomposed into sub-cubes of side $2^{-n}$, the inquiry whose explananda are the masses of the cubes at level $n$ treats level $n$ as fundamental, since those masses are fixed by the facts at level $n$, whatever the finer distribution. By Theorem~4.3 of MET-10-a, no level is fundamental absolutely, so any fundamentality a level has is of the relative kind.
\end{proof}

\section{Why the sciences are entitled to it}

The screening-off clause is what makes a level fit to be treated as fundamental. Chemistry may treat atoms as its basic units because the chemical properties it explains, valence and bonding behavior, are fixed by the electronic structure of atoms and do not vary with the internal states of their nuclei within the range that chemical processes reach. The special sciences are, on this account, autonomous in the sense of \citet{Fodor1974} without being independent of what lies below: their explananda are grounded in finer facts and screened off from them. The same pattern underlies the claim of MND-11 that mental causes are proportional and physical causes overly specific: the finer level carries more detail than the explanandum can register \citep{Yablo1992}.

Two consequences follow. First, a science that treats a level as fundamental, and then discovers a finer level, has not been refuted, provided clause (ii) still holds for its explananda; it has learned that its fundamentality was relative, which is the only kind there was. The picture of \citet{OppenheimPutnam1958}, on which the levels of the sciences rest on a bottom level of elementary particles, survives in every respect except the bottom \citep{Schaffer2003}. Second, the relativity is to the questions asked, not to anyone's interests in what exists. Which level is fundamental for an inquiry depends on the inquiry's explananda; what exists at each level, and what grounds what, do not depend on any inquiry at all.

Nor does the relative notion compete with the thesis of MET-09 that the cosmos is the one fundamental object. The cosmos is not a level, and its fundamentality is absolute: it depends on nothing. Levels, which are pluralities of proper parts, have only the fundamentality that inquiries confer on them by the questions they ask, and the descent ensures that no inquiry can find a level that does more.

\begin{thebibliography}{9}\small
\bibitem[Fodor(1974)]{Fodor1974} Fodor, J.~A. (1974). Special sciences (or: The disunity of science as a working hypothesis). \emph{Synthese}, 28(2), 97--115.
\bibitem[Oppenheim and Putnam(1958)]{OppenheimPutnam1958} Oppenheim, P. and Putnam, H. (1958). Unity of science as a working hypothesis. In H.~Feigl, M.~Scriven and G.~Maxwell (eds.), \emph{Minnesota Studies in the Philosophy of Science}, Volume 2. University of Minnesota Press.
\bibitem[Schaffer(2003)]{Schaffer2003} Schaffer, J. (2003). Is there a fundamental level? \emph{No\^us}, 37(3), 498--517.
\bibitem[Yablo(1992)]{Yablo1992} Yablo, S. (1992). Mental causation. \emph{Philosophical Review}, 101(2), 245--280.
\end{thebibliography}

\end{document}
